\documentclass[11pt]{article}
\usepackage{amsmath,amssymb,amsthm}
\usepackage[margin=1in]{geometry}
\newtheorem{theorem}{Theorem}
\newtheorem{lemma}{Lemma}
\begin{document}
\title{V112: Serial Phase Transfer for Ternary MUX Range Avoidance}
\author{NC0\_k-Avoid Laboratory}
\date{}
\maketitle

A signed essential ternary MUX gate has the form
\[
 g=o\oplus \operatorname{MUX}(s\oplus p_s,a\oplus p_a,b\oplus p_b).
\]
Fixing its output gives two exact binary clauses.  Branch $r$ supplies an implication
whose source phase is fixed by $p_s$ and whose arrival phase is selected by the output
target.

\paragraph{Serial two-lobe support.}
One root input selects two central outputs.  Each layer consists of two disjoint
two-variable lobes whose four selector gates return to a common exit variable.  A
non-final exit selects one shared MUX whose two data branches enter the next layer,
one branch per lobe.  The last layer exits to the root.  With $d$ shared hubs,
\[
 n=5(d+1),\qquad m=n+1.
\]
All prescribed gates may carry arbitrary signed MUX polarities.

\begin{lemma}[Recognition]
The exact serial two-lobe support template is recognizable in linear time.
\end{lemma}

\begin{proof}
The root is the unique selector of two outputs.  The two central gates expose four
first-layer variables.  Their unique selector supports share exactly one outside
variable, the first exit, and pair the four variables into two mutual lobes.  At an
internal exit $z$, its unique selector gate exposes one start variable from each next
lobe; the two start-selector supports intersect in the next exit and expose the two
partners.  Walking this pattern processes every gate and variable only constantly many
times, and the recognizer rejects unless the walk returns to the root using the entire
instance.
\end{proof}

For a shared hub traversal $(h,r)$ followed by first lobe branch $(g,c)$, define the
required target bit on $h$ by
\[
 t_h=\alpha(g,c)\oplus p_{\mathrm{data}}(h,r)\oplus o_h.
\]

\begin{lemma}[Local phase independence]
Suppose the two routes leave every shared hub into distinct lobes.  Whether their
target requirements agree at one hub is independent of every later layer.
\end{lemma}

\begin{proof}
The two lobes have disjoint output gates and meet only at their common exit.  Hence no
lobe gate introduces a cross-route target equality.  The current hub target depends
only on its selected branch and the source phase of the first selected branch in the
chosen lobe.  Both routes then reach the next common hub, so no other local information
is carried forward.
\end{proof}

\begin{theorem}[Serial phase-transfer avoider]
After recognition, there is a deterministic linear-time algorithm that decides whether
the circuit admits a lobe-disjoint phase-transfer certificate and constructs a missing
output whenever it does.
\end{theorem}

\begin{proof}
Enumerate the constant number of central branch pairs having opposite source phases and
entering distinct first lobes.  A two-variable lobe has only constantly many simple
branch paths to its exit.  At every shared hub enumerate the constant number of branch
pairs entering distinct next lobes and local lobe paths, retaining a choice exactly
when the two induced values of $t_h$ agree.  By local phase independence, successful
choices concatenate without global backtracking.

For the resulting cycles, target each selected gate so its implication arrives at the
source phase of the next branch, except the final branch, which arrives at the
complement of that cycle's initial root phase.  Shared hub targets agree by
construction.  Since the two central source phases are opposite, the two cycles force
opposite Boolean values on the root.  The fixed-output 2-CNF is therefore inconsistent,
and arbitrary values on unused outputs complete a missing word.
\end{proof}

\begin{theorem}[Minimum overlap]
Every pair of return routes in a depth-$d$ serial chain uses all $d$ shared hub outputs.
Lobe-disjoint routes therefore have minimum overlap exactly $d$.
\end{theorem}

\begin{proof}
Every layer exits at its hub.  Each non-final hub selects exactly one output gate, which
is the only way to leave that hub toward later layers.  Thus every return route uses
every shared hub.  Distinct lobes share no other output.
\end{proof}

\begin{theorem}[Mixed optimum face]
For every $d\ge1$ there is a periodic signed labeling of the serial support with both a
target-compatible and a target-incompatible pair at the same minimum overlap $d$.
\end{theorem}

\begin{proof}
Use the five-gate periodic signing recorded in the reference implementation.  Every
shared hub has data polarities $(1,1)$ and output flip zero.  For the compatible pair,
route zero takes hub branch one followed by the first right-lobe gate on branch one,
while route one takes hub branch zero followed by the first left-lobe gate on branch
zero.  The corresponding next source phases are respectively one and one, so both hub
requirements are zero.

For the incompatible pair, route zero takes hub branch zero followed by the first
left-lobe gate on branch one, whose source phase is zero, whereas route one takes hub
branch one followed by the first right-lobe gate on branch one, whose source phase is
one.  Their hub targets are therefore one and zero.  This repeats at every layer.  Both
pairs share exactly the mandatory $d$ hubs and no lobe output, so both are overlap
optimal.
\end{proof}

\begin{theorem}[Hall minimality]
The serial support is Hall-minimal for every depth: deleting any one output leaves a
perfect matching from the remaining $n$ outputs to the $n$ inputs.
\end{theorem}

\begin{proof}
Match one central output to the root and every non-central output to its selector,
leaving only the second central output unmatched.  If a central output is deleted, use
the other one at the root.  If a non-central selector output of input $x$ is deleted,
$x$ becomes free.  From the unmatched second central output, alternating support paths
reach every non-root input: its data variables enter both first lobes; matched lobe
gates reach their partners and layer exits; matched hub gates reach both starts of the
next layer.  Flip an alternating path ending at $x$ to obtain a perfect matching.
\end{proof}

\begin{theorem}[Exact affine backdoor]
On the periodic family,
\[
 \beta_{\mathrm{V102}}=3(d+1)=\frac{3n}{5}.
\]
\end{theorem}

\begin{proof}
The statement is support-only.  For each two-variable lobe, selecting its exit hub
requires at least one selected lobe variable; omitting the exit hub forces both lobe
variables.  Summing over two lobes in each of $d+1$ layers gives a lower bound
$3(d+1)$.  Selecting all $d+1$ exit hubs and one variable from every lobe attains it.
\end{proof}

\paragraph{Boundary.}
The algorithm is complete only for the stated lobe-disjoint serial certificate class.
It does not decide target compatibility over the full optimum-flow face of an arbitrary
MUX circuit, prove all essential MUX/bijunctive circuits are in P, solve unrestricted
$NC^0_3$-Avoid, establish a new general circuit lower bound, or resolve P versus NP.
Novelty and priority are not established.

\end{document}
